\documentclass[11pt]{article}
\usepackage[T1]{fontenc}
\usepackage{lmodern}
\usepackage[margin=1in]{geometry}
\usepackage{amsmath,amssymb,amsthm,mathtools}
\usepackage{microtype}
\usepackage{enumitem}
\usepackage{booktabs}
\usepackage[colorlinks=true,linkcolor=blue!45!black,citecolor=blue!45!black,urlcolor=blue!45!black]{hyperref}
\usepackage{xcolor}
\usepackage[nameinlink,noabbrev,capitalise]{cleveref}
\numberwithin{equation}{section}
\newtheorem{theorem}{Theorem}[section]
\newtheorem{lemma}[theorem]{Lemma}
\newtheorem{corollary}[theorem]{Corollary}
\newtheorem{conjecture}[theorem]{Conjecture}
\theoremstyle{definition}
\newtheorem{definition}[theorem]{Definition}
\newtheorem{remark}[theorem]{Remark}
\newtheorem{example}[theorem]{Example}
\newcommand{\pref}{\operatorname{pref}}
\newcommand{\suf}{\operatorname{suf}}
\newcommand{\len}[1]{\lvert #1\rvert}
\newcommand{\eps}{\varepsilon}
\newcommand{\ZZ}{\mathbb Z}
\newcommand{\Pre}{\mathcal P}
\newcommand{\Suf}{\mathcal Q}
\newcommand{\Ov}{\mathcal O}
\newcommand{\lsh}[1]{#1^{\leftarrow}}
\setlist{itemsep=3pt,topsep=5pt}
\title{The forced-count graph lives on prefixes and suffixes\\[2pt]
\large Notes toward a practical version of the polynomial-time SCS 2-approximation}
\date{October 8, 2026}
\author{}
\begin{document}
\maketitle

\begin{abstract}
We study the count construction of Section~2 of \emph{A Polynomial-Time
2-Approximation for Shortest Common Superstring} (OpenAI, September 2026),
which is defined over all $O(L^2)$ distinct substrings of the input. We prove
that its base graph is supported on the at most $2L$ prefixes and suffixes of
the input strings, and that the periodicity rule may be restricted, without
changing any count, to pairs of words that are each simultaneously a prefix and a suffix of
input strings. Both statements follow from a single \emph{shift lemma}. Finally,
we give a closed form for the blocking sums in terms of the edge multiplicities,
which yields an algorithm that never materializes the substring set. We also
show that restricting the rule to least periods is \emph{not} sound for this
purpose.
\end{abstract}

\section{Setting}

We use the notation of the preprint without change: the reduced instance
$\mathcal S$ (substring-free, no empty or duplicate strings), the set $V$ of
its distinct substrings, the blocking functional $B_A(s;m)$ of (2.2), the
counts $m$ defined by the extension rules (2.4), the requirement rule, and
the period rule (2.5), and the base multiplicities
$u(s)=m(s)-\sum_c m(cs)$, $d(s)=m(s)-\sum_c m(sc)$. Missing words have count
zero. Write
\[
\lambda(s)=\sum_{c\in\Sigma}m(cs),\qquad \rho(s)=\sum_{c\in\Sigma}m(sc),
\]
so that $u(s)=m(s)-\lambda(s)$ and $d(s)=m(s)-\rho(s)$. Let $\Pre$ and $\Suf$
be the sets of nonempty prefixes and nonempty suffixes of the strings in
$\mathcal S$, and let $\Ov=\Pre\cap\Suf$ be the set of \emph{overlap words}.
Clearly $\len\Pre,\len\Suf\le L$, and every required string lies in $\Ov$.

\begin{definition}
A \emph{rule} is a quadruple $(A,v,w,k)$ as in (2.5): $A$ is a periodic text
of period $p$, $v=A[0:\len v)\in V$, $w$ is a nonempty prefix of $v$, and
$\len v=\len w+kp$ with $k\ge1$. It is \emph{triggered} if
$m(v)>B_A(v;m)$; it then imposes $m(w)\ge B_A(w;m)+k+1$, its \emph{bound}.
By construction, $m(s)$ is the maximum of $\lambda(s)$, $\rho(s)$, the
requirement indicator, and the bounds of all triggered rules with $w=s$.
\end{definition}

For a text $A$ write $\lsh A$ for its translate $\lsh A(i)=A(i-1)$. If
$A[0:\len s)=s$ and $a=A(-1)$, then $\lsh A[0:\len s+1)=as$, so
$B_{\lsh A}(as;m)$ is defined. Translating a period-$p$ text yields a
period-$p$ text.

Two facts about $v$ and $w$ in a rule are used repeatedly. Since $w$ is a
prefix of $v$, $w\notin\Pre$ implies $v\notin\Pre$. Since $v$ has period $p$
and $\len v-\len w=kp$, also $w=v[kp:\len v)$, so $w$ is a suffix of $v$, and
$w\notin\Suf$ implies $v\notin\Suf$.

\section{Two elementary lemmas}

\begin{lemma}[Prefix-free sums]\label{lem:pfree}
Let $y$ be a nonempty word and $Z$ a finite prefix-free set of nonempty
words. Then $\sum_{z\in Z}m(yz)\le m(y)$. Symmetrically, if $Z$ is
suffix-free, then $\sum_{z\in Z}m(zy)\le m(y)$.
\end{lemma}

\begin{proof}
Induct on the maximum length in $Z$. Group $Z$ by first letter $c$. If the
one-letter word $c$ lies in $Z$, prefix-freeness makes it the only member
beginning with $c$. Otherwise $Z_c=\{z':cz'\in Z\}$ is prefix-free, and
induction gives $\sum_{z'\in Z_c}m(ycz')\le m(yc)$. Hence
$\sum_{z\in Z}m(yz)\le\sum_c m(yc)\le m(y)$ by the right extension rule.
(If $y\notin V$, every term is zero.) The suffix-free case is the mirror image,
using the left extension rule.
\end{proof}

\begin{lemma}[Bracket decomposition]\label{lem:decomp}
Let $A$ be periodic with $A[0:\len s)=s$, and put $a=A(-1)$, $e=A(\len s)$.
Then
\begin{align}
B_A(s;m)&=\Lambda_A(s)+B_{\lsh A}(as;m),&\Lambda_A(s)&\le\sum_{c\ne a}m(cs),\label{eq:left-decomp}\\
B_A(s;m)&=R_A(s)+B_A(se;m),&R_A(s)&\le\sum_{d\ne e}m(sd),\label{eq:right-decomp}
\end{align}
where $\Lambda_A(s)$ counts, with multiplicity $m(r)$, the bracketing
placements of $s$ whose left mismatch is immediately before $s$, and
$R_A(s)$ those whose right mismatch is immediately after $s$.
\end{lemma}

\begin{proof}
Consider a counted placement $j$ of $s$ in $r$ as in (2.1)--(2.2).

\emph{Left.} If $j\ge2$, then position $j-1$ is interior, so $r(j-1)=A(-1)=a$
and $as$ occurs at $j'=j-1\ge1$. Substituting $A(i)=\lsh A(i+1)$ turns each
condition of (2.1) for $(A,s,j)$ into the corresponding condition for
$(\lsh A,as,j')$: the length condition $j'+\len{as}=j+\len s\le\len r-1$, the
interior agreement $r(h)=A(h-j)=\lsh A(h-j')$, and the two end mismatches
$r(0)\ne A(-j)=\lsh A(-j')$ and
$r(\len r-1)\ne A(\len r-1-j)=\lsh A(\len r-1-j')$. The converse substitution
is equally direct, so these placements are in bijection with the counted
placements of $as$ relative to $\lsh A$, giving the second term. If $j=1$, then
$r=c\,s\,x\,d$ with $c\ne a$, $x=A[\len s:\len s+\len x)$ and
$d\ne A(\len s+\len x)$; the placement is unique for this $r$. For fixed $c$
the tails $x d$ form a prefix-free set, since two of them first differ where
one leaves $A$. \Cref{lem:pfree} bounds their total by $m(cs)$.

\emph{Right.} Mirror image: a placement whose right mismatch is not adjacent
to $s$ has $r(j+\len s)=A(\len s)=e$ and is exactly a counted placement of
$se$ relative to $A$. The adjacent ones have the form $r=y\,s\,d$ with
$d\ne e$, where the heads $y=c\,A[-\len y+1:0)$, $c\ne A(-\len y)$, form a
suffix-free set.
\end{proof}

\begin{corollary}[Monotonicity of the unblocked surplus]\label{cor:mono}
Put $U_A(s)=m(s)-B_A(s;m)$. Then $U_A(s)\ge0$, $U_A(s)\ge U_A(se)$ and
$U_A(s)\ge U_{\lsh A}(as)$. In particular, for fixed $A$ and $w$, the
rules $(A,A[0:\len w+kp),w,k)$ are triggered exactly for $k$ up to some
threshold, and only the largest triggered $k$ gives a binding bound.
\end{corollary}

\begin{proof}
By \eqref{eq:right-decomp} and the extension rule,
$m(s)\ge m(se)+\sum_{d\ne e}m(sd)\ge m(se)+R_A(s)$; subtract
$B_A(s;m)=R_A(s)+B_A(se;m)$. The left inequality is the mirror image.
Nonnegativity follows by induction on decreasing length, since $U_A(se)=0$
when $se\notin V$. If $A[0:\len w+kp)\in V$, so are its shorter prefixes, and
the bound $B_A(w;m)+k+1$ increases with $k$.
\end{proof}

\section{The shift lemma}

\begin{lemma}[Shift lemma]\label{lem:shift}
Let $(A,v,w,k)$ be a triggered rule.
\begin{enumerate}[label=\textup{(\alph*)}]
\item If $m(v)=\lambda(v)$, then $\lambda(w)\ge B_A(w;m)+k+1$.
\item If $m(v)=\rho(v)$, then $\rho(w)\ge B_A(w;m)+k+1$.
\end{enumerate}
\end{lemma}

\begin{proof}
(a) Let $a=A(-1)$; it is the left neighbor in $A$ of both $v$ and $w$. By
\eqref{eq:left-decomp} for $v$ and the trigger,
\[
m(av)+\sum_{c\ne a}m(cv)=\lambda(v)=m(v)>B_A(v;m)
=\Lambda_A(v)+B_{\lsh A}(av;m)\le\sum_{c\ne a}m(cv)+B_{\lsh A}(av;m).
\]
Hence $m(av)>B_{\lsh A}(av;m)\ge0$, so $av\in V$. Now
$(\lsh A,av,aw,k)$ is a rule: $\lsh A$ has period $p$,
$av=\lsh A[0:\len v+1)\in V$, $aw$ is a nonempty prefix of $av$, and
$\len{av}-\len{aw}=kp$. It is triggered, so the definition of $m$ at the
longer word $aw$ gives $m(aw)\ge B_{\lsh A}(aw;m)+k+1$. Applying
\eqref{eq:left-decomp} to $w$,
\[
\lambda(w)=m(aw)+\sum_{c\ne a}m(cw)\ge B_{\lsh A}(aw;m)+k+1+\Lambda_A(w)
=B_A(w;m)+k+1.
\]

(b) Let $e=A(\len v)$. Since $\len v-\len w=kp$, also $e=A(\len w)$. By
\eqref{eq:right-decomp} for $v$ and the trigger, exactly as in (a),
$m(ve)>B_A(ve;m)$, so $ve\in V$. Then $(A,ve,we,k)$ is a triggered rule:
$ve=A[0:\len v+1)$, its prefix $A[0:\len w+1)=we$, and
$\len{ve}-\len{we}=kp$. Hence $m(we)\ge B_A(we;m)+k+1$, and
\eqref{eq:right-decomp} for $w$ gives
$\rho(w)=m(we)+\sum_{d\ne e}m(wd)\ge B_A(we;m)+k+1+R_A(w)=B_A(w;m)+k+1$.
\end{proof}

\begin{remark}
In words: if every counted occurrence of $v$ extends to the left, then the
$k+1$ unblocked copies of $w$ that a rule extracts from $v$ can be shifted one
letter to the left, as copies of $aw$ in $av$. The rule is then already
implied by the left extension inequality at $w$. Part (b) uses the extension
to the right, together with the period: $w$ and $v$ have the same right
neighbor $e$ in $A$.
\end{remark}

\section{Main results}

\begin{theorem}[Prefix/suffix support]\label{thm:support}
For every nonempty $s\in V$:
\[
s\notin\Pre\implies m(s)=\lambda(s),\qquad
s\notin\Suf\implies m(s)=\rho(s).
\]
Consequently $u(s)>0$ implies $s\in\Pre$, and $d(s)>0$ implies
$s\in\Suf$. Every up edge of the base graph $G$ joins two words of
$\Pre\cup\{\eps\}$, every down edge joins two words of
$\Suf\cup\{\eps\}$, and $G$ has at most $\len{\Pre\cup\Suf}+1\le2L+1$
vertices.
\end{theorem}

\begin{proof}
Induct on decreasing length, proving both implications together. Let
$s\notin\Pre$. Then $s$ is not required. For each letter $d$ with $sd\in V$,
also $sd\notin\Pre$, so induction gives $m(sd)=\lambda(sd)$. Therefore
\[
\rho(s)=\sum_d\sum_c m(csd)=\sum_c\sum_d m(csd)\le\sum_c m(cs)=\lambda(s),
\]
by the right extension rule at each $cs$. Let $(A,v,s,k)$ be a triggered rule
with target $s$. Then $v\notin\Pre$ and $\len v>\len s$, so $m(v)=\lambda(v)$
by induction, and \cref{lem:shift}(a) bounds the rule by $\lambda(s)$. Every
lower bound defining $m(s)$ is therefore at most $\lambda(s)$, and
$m(s)=\lambda(s)$.

Let $s\notin\Suf$. Symmetrically, $\lambda(s)\le\rho(s)$, using induction on
$cs\notin\Suf$. For a triggered rule $(A,v,s,k)$, $s$ is a suffix of $v$, so
$v\notin\Suf$. Induction gives $m(v)=\rho(v)$, and \cref{lem:shift}(b)
bounds the rule by $\rho(s)$.

For the consequences, $u(s)=m(s)-\lambda(s)$ and $d(s)=m(s)-\rho(s)$. An up
edge $\pref(s)\to s$ is used only if $s\in\Pre$, and then
$\pref(s)\in\Pre\cup\{\eps\}$. Similarly, a down edge $s\to\suf(s)$ is used only
if $s\in\Suf$, and then $\suf(s)\in\Suf\cup\{\eps\}$.
\end{proof}

\begin{theorem}[Overlap-word restriction]\label{thm:restrict}
Let $m'$ be defined by the same descending recursion as $m$, but imposing
(2.5) only for rules with $v\in\Ov$ and $w\in\Ov$. Then $m'=m$ on all of $V$;
in particular, both define the same base graph.
\end{theorem}

\begin{proof}
Induct on decreasing length; assume $m'=m$ on all words longer than $s$.
The quantities $\lambda(s)$, $\rho(s)$, $B_A(s;\cdot)$, and the triggers of
all rules with target $s$ depend only on longer words, so they coincide for
$m$ and $m'$. Since $m'(s)$ is the maximum of a subset of the bounds defining
$m(s)$, we have $m'(s)\le m(s)$. Conversely, consider a triggered rule
$(A,v,s,k)$ omitted from the restricted system, so $s\notin\Ov$ or
$v\notin\Ov$. In either case $v\notin\Ov$, because $s\notin\Pre$ forces
$v\notin\Pre$ and $s\notin\Suf$ forces $v\notin\Suf$. If $v\notin\Pre$, \cref{thm:support} gives $m(v)=\lambda(v)$ and
\cref{lem:shift}(a) bounds the rule by $\lambda(s)\le m'(s)$. If
$v\notin\Suf$, part (b) bounds it by $\rho(s)\le m'(s)$. Thus
$m(s)\le m'(s)$.
\end{proof}

\begin{corollary}\label{cor:determined}
Every count is determined by the counts on $\Pre\cup\Suf$: for
$s\notin\Pre$, $m(s)=\sum_c m(cs)$, and for $s\notin\Suf$,
$m(s)=\sum_c m(sc)$. The only period rules that must be evaluated are those with
$v,w\in\Ov$, and, by \cref{cor:mono}, for each overlap word $w$ and each
compatible periodic text $A$ only the largest triggered $k$.
\end{corollary}

\section{Least periods do not suffice}

The rule (2.5) ranges over all periods of $v$, not only its least period. Restricting the rule to least
periods weakens the bound, even on overlap words.

\begin{example}\label{ex:least}
Let $\mathcal S=\{\texttt{aabaa},\texttt{baaab}\}$, with optimum $7$
(\texttt{baaabaa}). The word $v=\texttt{aabaa}\in\Ov$ has least period $3$ and
also period $4$. With $A=(\texttt{aaba})^\infty$, the rule $(A,v,\texttt{a},1)$
is triggered and gives $m(\texttt{a})\ge B_A(\texttt{a};m)+2=2+2=4$. All other
bounds on $m(\texttt{a})$ are at most $3$. With all periods, $W=6$; with least
periods only, $m(\texttt{a})=3$ and $W=5$. The base-graph lower bound is
therefore genuinely weaker under the least-period restriction, and any
efficient implementation must handle all periods of overlap words.
\end{example}

\section{Empirical confirmation and what remains}

A direct implementation of the preprint's Section~2 confirmed every statement above,
with no violations: \cref{thm:support} and both parts of \cref{lem:shift}
(about $5\times10^6$ triggered rules), \cref{cor:mono}, and
\cref{thm:restrict} (identical base graphs on $4\times10^5$ random,
periodic, and Fibonacci-type instances).

The following observation is not proved here.

\begin{conjecture}
Without the period rule, $W=\sum_c m(c)$ equals the weight of a minimum cycle
cover of the overlap (distance) graph of $\mathcal S$.
\end{conjecture}

It held on all $1500$ instances tested. If true, the extension-only part of
the construction is computable in $O(L)$ time by known Aho--Corasick
techniques, and only the periodic correction is new.

\paragraph{Consequence for the algorithm.}
\Cref{thm:support,thm:restrict} reduce the vertex set from $O(L^2)$ to
$\len{\Pre\cup\Suf}\le2L$, and the period rules to pairs of overlap words. On
600 reads of length 100 from a random genome this is $1.25\times10^6$
substrings versus $1.1\times10^5$ prefixes and suffixes and $2{,}879$ overlap
words, with the number of $(v,p)$ rule sources falling from about
$4\times10^5$ to $944$. The blocking sums are handled in the next section.

\section{Blocking sums from the edge multiplicities}\label{sec:blocking}

The blocking functional $B_A(w;m)$ ranges over bracketing words that need not
lie in $\Pre\cup\Suf$. We now express it, and the trigger of a rule, using
only $u$ and $d$, which by \cref{thm:support} live on $\Pre$ and $\Suf$
respectively and have total multiplicity $W$ each.

Throughout, $A$ is a periodic text with $A[0:\len x)=x$, and we write
$x_n=A[0:n)$ for $n\ge\len x$.

\begin{lemma}[Left-bracket telescoping]\label{lem:left-tele}
For a nonempty word $y$ whose left context is read in $A$ (that is, we put
$y_\mu=A[-\mu:0)\,y$), define
\[
\Lambda^*_A(y)=\sum_{\mu\ge0}\ \sum_{c\ne A(-\mu-1)} m(c\,y_\mu),
\]
the total count of words obtained from $y$ by extending along $A$ to the left
and then adding one mismatching letter. Then
$\Lambda^*_A(y)=m(y)-\sum_{\mu\ge0}u(y_\mu)$.
\end{lemma}

\begin{proof}
For each $\mu$, $\sum_c m(c\,y_\mu)=\lambda(y_\mu)=m(y_\mu)-u(y_\mu)$, and the
term with $c=A(-\mu-1)$ is $m(y_{\mu+1})$. Hence the $\mu$-th inner sum is
$m(y_\mu)-u(y_\mu)-m(y_{\mu+1})$. Summing over $\mu$ telescopes. The sum is
finite, since $y_\mu\notin V$ for large $\mu$, and then all terms vanish.
\end{proof}

\begin{theorem}[Closed form for blocking sums]\label{thm:closed}
Put
\[
D_A(x)=\sum_{n\ge\len x}d(x_n),\qquad
T_A(x)=\sum_{n\ge\len x}\ \sum_{c\ne A(n)}\ \sum_{\mu\ge0}
u\bigl(A[-\mu:n)\,c\bigr).
\]
Then the unblocked surplus is
\begin{equation}\label{eq:U-closed}
U_A(x)=m(x)-B_A(x;m)=D_A(x)+T_A(x),
\end{equation}
and equivalently
\begin{equation}\label{eq:B-closed}
B_A(x;m)=\rho(x)-\bigl(D_A(x)-d(x)\bigr)-T_A(x).
\end{equation}
\end{theorem}

\begin{proof}
In \cref{lem:decomp}, the bracketing placements of $x_n$ whose right mismatch
is adjacent are exactly the words $c'\,A[-\mu:0)\,x_n\,c$ with
$c'\ne A(-\mu-1)$ and $c\ne A(n)$, each with one placement. Hence, with
equality, $R_A(x_n)=\sum_{c\ne A(n)}\Lambda^*_A(x_n c)$. Applying
\cref{lem:left-tele} to $y=x_nc$, and using
$\sum_{c\ne A(n)}m(x_nc)=\rho(x_n)-m(x_{n+1})=m(x_n)-d(x_n)-m(x_{n+1})$,
\[
R_A(x_n)=m(x_n)-m(x_{n+1})-d(x_n)
-\sum_{c\ne A(n)}\sum_{\mu\ge0}u\bigl(A[-\mu:n)\,c\bigr).
\]
Iterating the right decomposition $B_A(x_n)=R_A(x_n)+B_A(x_{n+1})$ from
$n=\len x$ until $x_n\notin V$, where everything vanishes, gives
$B_A(x;m)=m(x)-D_A(x)-T_A(x)$. Finally $m(x)-d(x)=\rho(x)$.
\end{proof}

\begin{remark}[Walk interpretation]
On the closed walks of Sections~2--3 of the preprint, \eqref{eq:U-closed} says
that an occurrence of $x$ is unblocked exactly when, before any mismatch has entered on both
sides, either the walk deletes the first letter of the occurrence while its
window is still an $A$-aligned word $x_n$ (the term $D_A$), or the walk
appends a first right mismatch $c$ while its window is an $A$-aligned word
covering $x$ (the term $T_A$). This matches the proof of Lemma~3.4 of the
preprint.
\end{remark}

\begin{corollary}[Counts from the supports of $u$ and $d$]\label{cor:uv-only}
Process the words of $\Pre\cup\Suf$ in decreasing length. For each such word $s$,
\[
\lambda(s)=\sum_{\substack{y\text{ longer, ends with }s}}u(y),\qquad
\rho(s)=\sum_{\substack{y\text{ longer, starts with }s}}d(y),
\]
by the telescoping identity in the proof of Lemma~3.1 of the preprint. Set $m(s)=\rho(s)$ if
$s\notin\Suf$ and $m(s)=\lambda(s)$ if $s\notin\Pre$. If $s\in\Ov$, take the maximum of
$\lambda(s)$, $\rho(s)$, the requirement indicator, and, for each pending rule $(A,k)$,
\[
\rho(s)+k+1-\bigl(D_A(s)-d(s)\bigr)-T_A(s).
\]
After fixing $m(s)$, set $u(s)=m(s)-\lambda(s)$ and $d(s)=m(s)-\rho(s)$. Then, for
$s\in\Ov$ and each period $p$ of $s$, the rule with $A$ generated by
$s[0:p)$ is triggered iff $D_A(s)+T_A(s)>0$. Every quantity used involves
$u$ and $d$ only on words already processed, together with $d(s)$ itself in
the trigger.
\end{corollary}

In particular a rule can raise $m(w)$ above $\rho(w)$ only if
$k+1>(D_A(w)-d(w))+T_A(w)$.

\begin{lemma}[Primitive templates suffice]\label{lem:primitive}
If $v[0:p)=t^j$ with $j\ge2$, every rule $(A,v,w,k)$ with period $p$ is
dominated by the rule $(A,v,w,jk)$ with period $\len t$.
\end{lemma}

\begin{proof}
Both rules use the same text $A=t^\infty$, hence the same trigger and the
same $B_A(w;m)$. Since $\len v-\len w=kp=(jk)\len t$, the second rule
exists, and its bound $B_A(w;m)+jk+1$ is larger.
\end{proof}

\paragraph{Implementation and validation.}
A prototype (\texttt{compact.py}) keeps $u$ in a forward and a reversed weighted
trie and $d$ in a forward trie, built incrementally in decreasing length.
Then $\lambda$ and $\rho$ are subtree sums, and $D_A$ is a single walk along $A$.
For $T_A$, words $A[-\mu:n)c$ with the same start residue modulo $p$ are the same
string, so $T_A(x)$ takes $p$ walks, each weighted by the number of admissible
shifts. The prototype never builds $V$. On $4.6\times10^5$ random, periodic
and Fibonacci-type instances it reproduced the reference implementation's $u$
and $d$ exactly. It processes $2\times10^4$ reads of length $100$
($L=2\times10^6$) in under 90\,s in pure Python, whereas the reference
construction would have to enumerate about $2\times10^7$ distinct substrings there and evaluate blocking sums over them. The remaining costs are $O(\len s)$ trie navigation per
word and the $T_A$ walks. Both should become amortized constant time per
step using Aho--Corasick node identifiers and Fenwick trees over Euler tours.

\end{document}
